\section{Proof of Theorem~\ref{thm-frame-bound}: the John-frame certificate bound}\label{sec:8}

We now prove Theorem~\ref{thm-frame-bound}, deferred in Section~\ref{sec:3.3}. After normalization,
its determinant assertion is Theorem~\ref{thm-determinant} below. Section~\ref{sec:8.6} will return to
the original certificate and complete both statements.

First normalize the determinant at the equality model. Section~\ref{sec:8.1} records
the scalar moments needed for the high-dimensional proof in Section~\ref{sec:8.2}.
Only after that proof do we introduce transverse blocks: Section~\ref{sec:8.3} gives
their basic constraints and settles dimension two. Section~\ref{sec:8.4} states and
proves the direct coupling estimate used in dimension three and local
stability, followed by the three-dimensional proof in Section~\ref{sec:8.5}.
Section~\ref{sec:8.6} combines the branches and recovers equality.

\textbf{From the certificate to a symmetric matrix.} Recall $b=1/n$,
$P=I-vv^T$, and $u_i=s_iv+y_i$, where $s_i=u_i^Tv$ and $y_i=Pu_i$.
Use the parameters of Section~\ref{sec:3}:
$d=2-b$, $a=b/d$, $t=(1+b)/2$, $k=2\sqrt{1-b}/d$, and
$f(s)=(1-ts)/(1-as)$. These parameters depend only on the dimension; the
contact data enter through the weighted second-moment matrix $B$ below.
The useful factorization of the existing certificate is
\begin{equation}\label{eq:8.1}
\begin{aligned}
B&=\sum_i\lambda_i f(s_i)u_iu_i^T,&
F&=B+\frac{t}{1-a}vv^T,\\
D&=kP+k^2vv^T,&Q&=FD.
\end{aligned}
\end{equation}

Here $F$ adds a rank-one axial term to $B$, while $D$ rescales the transverse
directions by $k$ and the axial direction by $k^2$. This separates the
frame-dependent symmetric matrix $F$ from a fixed scaling whose determinant
is known.

To check the factorization, use
$\sum_i\lambda_i f(s_i)u_i=aBv-tv$ in the expansion of \eqref{eq:3.11}.
It gives $Q=BD+t(1+a)vv^T=FD$. All $f(s_i)$ are positive, so isotropy
makes $B$ and $F$ positive definite. Also $D\succ0$. This already proves
that $Q$ is invertible and has positive determinant.

\textbf{Choose the equality model as the reference.} At the model contacts, the
certificate is $Q=M_0=A_*^{-1}$. Thus the corresponding symmetric factor is
$F_0=M_0D^{-1}$. Normalize $F$ by this reference matrix:
\begin{equation}\label{eq:8.2}
\begin{gathered}
F_0=\frac{d^2}{2(1-b)}vv^T+\frac d2P,
\qquad\widehat F=F_0^{-1/2}FF_0^{-1/2},\\
F_0D=M_0,\qquad\det\widehat F=R_n\det Q.
\end{gathered}
\end{equation}

This congruence preserves positive definiteness and sends $F_0$ to the identity.
Consequently, $\widehat F-I$ measures the matrix departure from the equality
model. The last identity uses $\det D=k^{n+1}$ and
$\det(F_0D)=\det M_0=R_n^{-1}$. Thus the normalized determinant theorem
below is exactly the quantitative assertion needed in Theorem~\ref{thm-frame-bound}.

To measure the same departure directly in the contact data, recall the scalar
defect from \eqref{eq:3.10}:
\begin{equation}\label{eq:8.3}
\chi=\sum_i\lambda_i\frac{(1+s_i)(s_i-b)^2}{1-as_i}\ge0.
\end{equation}

The factors $1+s_i$ and $(s_i-b)^2$ make each summand vanish at the model pole
and latitude, respectively. The denominator is positive, and the John weights
retain the contribution of each contact. The next theorem links this scalar
description of contact geometry to the normalized determinant.

\begin{theorem}[{quantitative John-frame determinant inequality}]\label{thm-determinant}
Every John frame in dimension $n\ge2$ satisfies
\begin{equation}\label{eq:8.4}
\frac{\chi}{1500}\le\log\det\widehat F\le\frac{\chi}{2d},
\qquad
\det F\ge\det F_0=\frac{d^{n+1}}{2^n(1-b)}.
\end{equation}

Equality in the determinant inequality holds if and only if $s_i\in\{-1,b\}$ for every contact, in which case $F=F_0$ and $\widehat F=I$.

\end{theorem}

\textbf{Proof developed in Sections~\ref{sec:8.1}--\ref{sec:8.6}.} The dimension-specific lower
bounds are established in Sections~\ref{sec:8.2}, \ref{sec:8.3}, and 8.5; Section~\ref{sec:8.6} assembles
them and proves the upper bound and equality statement.

The upper estimate will follow from the trace of $\widehat F-I$. The main work
is the lower estimate: a positive trace alone does not give the required determinant lower bound,
so we must use the relations between the matrix entries imposed by the contacts.
The coefficient $1/1500$ is the explicit margin produced by these estimates.
The sharp geometric constant is $R_n$; the proof does not assert optimality
of this quantitative margin.

\subsection{Scalar preparation for Theorem~\ref{thm-determinant}}\label{sec:8.1}

For the high-dimensional argument, the needed contact information is the
height $s_i$ and the scalar defect $\chi$. Reweight by
$\nu_i=\lambda_i/(1-as_i)$ to absorb the denominator in $f$, and write
$N_r=\sum_i\nu_i s_i^r$. The scalar John moments give
\begin{equation}\label{eq:8.5}
N_0-aN_1=n,\qquad N_1-aN_2=0,\qquad N_2-aN_3=1.
\end{equation}

The three relations express $N_0,N_1,N_3$ in terms of $N_2$.
Expanding the defect gives
\begin{equation}\label{eq:8.6}
\chi=N_3+(1-2b)N_2+(b^2-2b)N_1+b^2N_0.
\end{equation}

Substitution expresses the axial entry of $B$ as
\begin{equation}\label{eq:8.7}
\zeta:=v^TBv=N_2-tN_3
=\frac{3d}{4}-\frac{d^2}{4(1-b)(2+b)}\chi.
\end{equation}

This is the scalar input used below: departure from the model lowers
$\zeta$ by a known multiple of $\chi$. All denominators are positive since
$0<b\le1/2$ and $-1\le s_i\le1$. The transverse information will be introduced
when the low-dimensional argument needs it.

\Needspace{22\baselineskip}

\subsection{Higher dimensions in Theorem~\ref{thm-determinant}}\label{sec:8.2}

We now prove the lower bound in Theorem~\ref{thm-determinant} for $n\ge4$, so $0<b\le1/4$.
Two lemmas isolate the quantities to compare: Lemma~\ref{lem-lifted-bound} gives a determinant
gain and leaves a scalar logarithmic difference; Lemma~\ref{lem-hermite} bounds that
difference. Lemma~\ref{lem-rational} then completes the constant comparison.

\subsubsection{The lifted determinant bound}\label{sec:8.2.1}

Set $\ell(s)=\log f(s)$ and $\widetilde u_i=(u_i,\sqrt b)$.
Define the lifted moment matrix and weighted logarithmic sum by
\begin{equation}\label{eq:8.8}
\mathcal A=\sum_i\lambda_i f(s_i)\widetilde u_i\widetilde u_i^T,
\qquad \mathcal L=\sum_i\lambda_i\log f(s_i).
\end{equation}

The corresponding value for the pole-and-latitude model is
\begin{equation}\label{eq:8.9}
\mathcal L_0:=\frac{\ell(-1)}{1+b}+\frac{\ell(b)}{b(1+b)}.
\end{equation}

The two coefficients are the total John weights at the model pole and latitude.
Define also the gain coefficient
\begin{equation}\label{eq:8.10}
\delta_0(b):=\frac{(1+b)(b^2+b+2)}
{2d(1-b)(2+b)(3+b)}.
\end{equation}

\begin{lemma}[{lifted determinant bound}]\label{lem-lifted-bound}
For every John frame in dimension
$n\ge4$,
\begin{equation}\label{eq:8.11}
\log\det\widehat F
\ge (1+b)(\mathcal L-\mathcal L_0)+\delta_0(b)\chi.
\end{equation}

\end{lemma}

The matrix problem is thus reduced to a lower bound on
$\mathcal L-\mathcal L_0$. The next subsection estimates precisely this term.

\begin{proof}
\textbf{Encode the moments in the lifted matrix.} The constant last
coordinate turns the zero first moment into zero off-diagonal blocks, and
$b=1/n$ makes the final diagonal entry equal to one.

John's first and second moments combine into $\sum_i\lambda_i\widetilde u_i\widetilde u_i^T=I_{n+1}$. We use the standard finite-frame determinant estimate
\begin{equation}\label{eq:8.12}
XX^T=I_r,\quad \rho_i>0
\quad\Longrightarrow\quad
\log\det\bigl(X\operatorname{diag}(\rho_i)X^T\bigr)
\ge\sum_i\|X_i\|^2\log\rho_i,
\end{equation}
where $X_i$ are the columns. The point of this estimate is to replace a matrix
determinant by a weighted scalar sum. In the standard Cauchy--Binet/Jensen
argument \cite{HJ1985,BV2004}, the coefficients $w_S=(\det X_S)^2$, for subsets
$S$ of $r$ column indices, form a probability distribution. The needed
determinant identities are
\begin{equation}\label{eq:8.13}
\sum_S w_S=1,\qquad
\det(X\operatorname{diag}(\rho_i)X^T)=\sum_Sw_S\prod_{i\in S}\rho_i,
\qquad \sum_{S\ni i}w_S=\|X_i\|^2.
\end{equation}

For the marginal identity, expand the same determinant when only the $i$th diagonal weight is $1+z$:
$\det(I_r+zX_iX_i^T)=1+z\|X_i\|^2$.
Comparing coefficients of $z$ suffices. Scalar Jensen now gives
\begin{equation}\label{eq:8.14}
\log\!\left(\sum_Sw_S\prod_{i\in S}\rho_i\right)
\ge\sum_Sw_S\sum_{i\in S}\log\rho_i
=\sum_i\|X_i\|^2\log\rho_i,
\end{equation}
which proves \eqref{eq:8.12}. For $X_i=\sqrt{\lambda_i}\widetilde u_i$, this becomes
\begin{equation}\label{eq:8.15}
\log\det\mathcal A\ge(1+b)\mathcal L.
\end{equation}

\textbf{Return to the original determinant.} Both $\mathcal A$ and $F$ contain
$B$, so their determinant ratio can cancel this common matrix factor.
Write $\kappa=t/(1-a)$ for the axial coefficient already present in $F$,
and $\varrho=v^TB^{-1}v$ for the inverse axial moment introduced by the
Schur complement. The same moment identities used in \eqref{eq:8.1} give
\begin{equation}\label{eq:8.16}
\mathcal A=
\begin{pmatrix}
B&\sqrt b(aBv-tv)\\
\sqrt b(aBv-tv)^T&1-abt+a^2b\zeta
\end{pmatrix},
\qquad
\det\mathcal A=\frac{\det B}{n}(n+at-t^2\varrho).
\end{equation}

The Schur complement is positive. The abbreviation $H_0=n+at$ records its
constant part. Our aim is to remove $\varrho$ in favor of the direct axial
moment $\zeta=v^TBv$, whose dependence on $\chi$ is known from \eqref{eq:8.7}.
The determinant lemma and Cauchy--Schwarz give
\begin{equation}\label{eq:8.17}
\det Q=k^{n+1}\det B(1+\kappa\varrho),\qquad
\varrho\zeta\ge1,\qquad H_0-t^2\varrho>0.
\end{equation}

Cross multiplication, with positive denominators, therefore yields
\begin{equation}\label{eq:8.18}
\frac{\det Q}{\det\mathcal A}
=nk^{n+1}\frac{1+\kappa\varrho}{H_0-t^2\varrho}
\ge nk^{n+1}\frac{\zeta+\kappa}{H_0\zeta-t^2}.
\end{equation}

Taking logarithms and applying \eqref{eq:8.15}, we obtain
\begin{equation}\label{eq:8.19}
\log\det Q\ge(n+1)\log k+(1+b)\mathcal L+g(\zeta),
\quad
g(\zeta)=\log\frac{n(\zeta+\kappa)}{(n+at)\zeta-t^2}.
\end{equation}

All quantities inside the logarithm are positive: \eqref{eq:8.16} gives
$n+at>t^2\varrho\ge t^2/\zeta$.

\textbf{Extract the positive gain from the axial correction.} As $\chi$ increases,
\eqref{eq:8.7} makes $\zeta$ decrease from its model value $\zeta_0=3d/4$.
The variables $\theta_1,\theta_2$ below measure the relative decreases in the
numerator and denominator inside $g$. The denominator decreases more, which
increases their ratio. The elementary inequality $\log x\ge1-1/x$ will quantify
this increase. Write
\begin{equation}\label{eq:8.20}
\Delta_\zeta=\zeta_0-\zeta
=\frac{d^2\chi}{4(1-b)(2+b)},\qquad
\theta_1=\frac{\Delta_\zeta}{\zeta_0+\kappa},\qquad
\theta_2=\frac{H_0\Delta_\zeta}{H_0\zeta_0-t^2}.
\end{equation}

The positive-denominator condition implies $0\le\theta_1\le\theta_2<1$. Hence
\begin{equation}\label{eq:8.21}
\begin{aligned}
g(\zeta)-g(\zeta_0)
&=\log\frac{1-\theta_1}{1-\theta_2}
\ge1-\frac{1-\theta_2}{1-\theta_1}\\
&=\frac{\theta_2-\theta_1}{1-\theta_1}
\ge\theta_2-\theta_1=\delta_0(b)\chi,
\end{aligned}
\end{equation}

The last equality follows by substituting the parameters into the definition
of $\delta_0(b)$.

\textbf{Subtract the model value.} Since $f(-1)=d(3+b)/4$ and $f(b)=d/2$,
the model sum defined in \eqref{eq:8.9} satisfies
\begin{equation}\label{eq:8.22}
(n+1)\log k+(1+b)\mathcal L_0+g(\zeta_0)=\log R_n^{-1}.
\end{equation}

Combine \eqref{eq:8.19}, \eqref{eq:8.21}, and \eqref{eq:8.22}, and use
$\log\det\widehat F=\log R_n+\log\det Q$ from \eqref{eq:8.2}.
This gives \eqref{eq:8.11}, completing the proof. 
\end{proof}

\Needspace{12\baselineskip}

\subsubsection{The interpolation loss}\label{sec:8.2.2}

Keep $\ell(s)=\log f(s)$ and $\mathcal L=\sum_i\lambda_i\ell(s_i)$.
Recall explicitly the model value
$\mathcal L_0=\ell(-1)/(1+b)+\ell(b)/[b(1+b)]$ from \eqref{eq:8.9}.
Define the scalar integral
\begin{equation}\label{eq:8.23}
I(b):=\int_a^t\frac{z^2}{(1-z^2)(1-bz)^2}\,dz.
\end{equation}

\begin{lemma}[{interpolation loss}]\label{lem-hermite}
With these definitions, every John frame
in dimension $n\ge4$ satisfies
\begin{equation}\label{eq:8.24}
\mathcal L-\mathcal L_0\ge-(1-a)I(b)\chi.
\end{equation}

\end{lemma}

This is the scalar estimate needed in Lemma~\ref{lem-lifted-bound}. Its proof also gives the
exact pointwise remainder \eqref{eq:8.27}.

\begin{proof}
\textbf{Choose a quadratic with the correct weighted sum.} The John identities determine
the weighted average of every quadratic polynomial in the heights. We therefore
replace $\ell(s)=\log f(s)$ by a quadratic that matches the equality heights.
Matching the value at $-1$ and both the value and slope at $b$ also gives the
same vanishing factors as the defect $\chi$. Let $H_2$ satisfy
$H_2(-1)=\ell(-1)$, $H_2(b)=\ell(b)$, and $H_2'(b)=\ell'(b)$. Explicitly,
\begin{equation}\label{eq:8.25}
H_2(s)=\ell(b)+\ell'(b)(s-b)
+\frac{\ell(-1)-\ell(b)+(1+b)\ell'(b)}{(1+b)^2}(s-b)^2.
\end{equation}

The scalar John moments agree with those of masses $1/(1+b)$ at $-1$ and $1/[b(1+b)]$ at $b$. Since $H_2$ is quadratic,
\begin{equation}\label{eq:8.26}
\sum_i\lambda_iH_2(s_i)=\mathcal L_0.
\end{equation}

Therefore $\mathcal L-\mathcal L_0$ is the weighted sum of the remainders
$\ell(s_i)-H_2(s_i)$. We will establish the exact identity, for $-1\le s\le1$,
\begin{equation}\label{eq:8.27}
\ell(s)-H_2(s)=-(1+s)(s-b)^2
\int_a^t\frac{z^2\,dz}{(1+z)(1-bz)^2(1-zs)}.
\end{equation}

\textbf{Derive and bound the remainder.} We first prove a rational interpolation identity, then integrate
it to recover the logarithm. This exposes the sign and the vanishing factors
of the remainder explicitly.

\textbf{Interpolate the rational integrand.} Put $r_z(s)=s/(1-zs)$ and let
$H_{2,z}$ be the quadratic formula \eqref{eq:8.25} with data $r_z(-1)$, $r_z(b)$
and slope $(1-zb)^{-2}$. Clearing denominators gives the exact rational identity
\begin{equation}\label{eq:8.28}
r_z(s)-H_{2,z}(s)
=\frac{z^2(1+s)(s-b)^2}{(1+z)(1-bz)^2(1-zs)}.
\end{equation}

\textbf{Recover the logarithmic interpolant.} The value data integrate to
$\ell(-1)$ and $\ell(b)$ because $\ell(s)=-\int_a^t r_z(s)\,dz$.
It remains to check the slope data; these satisfy the elementary identity
\begin{equation}\label{eq:8.29}
-\int_a^t\frac{dz}{(1-bz)^2}
=-\frac{t-a}{(1-ab)(1-tb)}
=\frac{a}{1-ab}-\frac{t}{1-tb}=\ell'(b).
\end{equation}

Integration of the three polynomial coefficients therefore gives \eqref{eq:8.27}, with
no differentiation under the integral. All integrands are continuous and all
denominators are positive on the compact interval.

\textbf{Express the loss in units of the contact defect.} To sum the remainder,
we need a coefficient independent of the individual height $s$. The following
comparison leaves the factor defining $\chi$ outside the integral:
\begin{equation}\label{eq:8.30}
\frac{1-a}{1-z}-\frac{1-as}{1-zs}
=\frac{(z-a)(1-s)}{(1-z)(1-zs)}\ge0
\qquad(a\le z\le t,\ -1\le s\le1).
\end{equation}

It follows pointwise that
$\ell(s)-H_2(s)\ge-(1-a)I(b)(1+s)(s-b)^2/(1-as)$.
Summing with weights $\lambda_i$ proves \eqref{eq:8.24}. 
\end{proof}

\textbf{Combine the gain and loss estimates.} Lemmas~\ref{lem-lifted-bound} and \ref{lem-hermite} now give the
two inputs to the higher-dimensional part of Theorem~\ref{thm-determinant}. Substituting the
interpolation bound into \eqref{eq:8.11} yields
\begin{equation}\label{eq:8.31}
\log\det\widehat F\ge\delta_0(b)[1-\mathcal R(b)]\chi,
\qquad
\mathcal R(b)=\frac{4(1-b)^2(2+b)(3+b)}{b^2+b+2}I(b).
\end{equation}

Here $\mathcal R(b)=(1+b)(1-a)I(b)/\delta_0(b)$ is the loss as a fraction of
the gain. The remaining task is the scalar comparison in Lemma~\ref{lem-rational}.

\subsubsection{Exact evaluation and the scalar certificate}\label{sec:8.2.3}

The matrix argument is now complete except for a scalar comparison:
$\mathcal R(b)$ measures the interpolation loss as a fraction of the gain
$\delta_0(b)\chi$. We need a uniform margin below one and a positive lower bound
for $\delta_0$. The convenient targets below give
$(1/12)(1-124/125)=1/1500$, the coefficient in Theorem~\ref{thm-determinant}.

\begin{lemma}[{uniform scalar certificate}]\label{lem-rational}
For $0<b\le1/4$,
$\mathcal R(b)<124/125$ and $\delta_0(b)\ge1/12$.

\end{lemma}

\begin{proof}
We first reduce to the endpoint $b=1/4$, evaluate the integral by
partial fractions, and compare its exact value with $124/125$.
The lower bound for $\delta_0$ is a separate polynomial comparison at the end.

\textbf{Step 1: reduce to the largest parameter.} In the original integral, both
endpoints and the integrand depend on $b$. The substitution
$w=(z-a)/(1-az)$ fixes the lower endpoint at zero and exposes three factors
that can each be compared as $b$ varies:
\begin{equation}\label{eq:8.32}
\mathcal R(b)=\frac{4(2+b)(3+b)}{b^2+b+2}
\int_0^{(2+b)/(4+b)}\frac1{1-w^2}
\left(\frac{2w+b(1-w)}{2+b(1-w)}\right)^2dw.
\end{equation}

Temporarily name the outside prefactor $\pi$, the rational factor inside the
integrand $\psi_b$, and the upper endpoint $\ell_b$:
$\pi(b)=4(2+b)(3+b)/(b^2+b+2)$,
$\psi_b(w)=[2w+b(1-w)]/[2+b(1-w)]$, and
$\ell_b=(2+b)/(4+b)$. The largest endpoint is $\ell_{1/4}=9/17$.
For $0\le b_1\le b_2\le1/4$ and $0\le w\le9/17$,
\begin{equation}\label{eq:8.33}
\begin{aligned}
\pi(b_2)-\pi(b_1)
&=\frac{16(b_2-b_1)(1-b_1-b_2-b_1b_2)}
{(b_1^2+b_1+2)(b_2^2+b_2+2)}\ge0,\\
\psi_{b_2}(w)-\psi_{b_1}(w)
&=\frac{2(b_2-b_1)(1-w)^2}
{[2+b_1(1-w)][2+b_2(1-w)]}\ge0,\\
\ell_{b_2}-\ell_{b_1}
&=\frac{2(b_2-b_1)}{(4+b_1)(4+b_2)}\ge0.
\end{aligned}
\end{equation}

Here $1-b_1-b_2-b_1b_2\ge7/16$ and $\psi_b(w)\ge0$. The integrand and interval in \eqref{eq:8.32} are therefore increasing in $b$, giving
$\mathcal R(b)\le\mathcal R(1/4)$.

\textbf{Step 2: evaluate the endpoint integral by partial fractions.}
Specializing the same $w$-integral \eqref{eq:8.32} to $b=1/4$ gives
\begin{equation}\label{eq:8.34}
\mathcal R(1/4)=\frac{468}{37}
\int_0^{9/17}\frac{(1+7w)^2}{(1-w^2)(9-w)^2}\,dw.
\end{equation}

Its integrand has the elementary decomposition
\begin{equation}\label{eq:8.35}
\begin{aligned}
\frac{(1+7w)^2}{(1-w^2)(9-w)^2}
&=\frac{9}{50(1+w)}+\frac{1}{2(1-w)}\\
&\quad-\frac{8}{25(9-w)}-\frac{256}{5(9-w)^2}.
\end{aligned}
\end{equation}

All logarithm arguments are positive on $0\le w\le9/17$.
Integrating the four terms and collecting the endpoint logarithms yields
\begin{equation}\label{eq:8.36}
\begin{aligned}
\mathcal R(1/4)
&=\frac{468}{37}
\left(\frac9{50}\log13-\frac1{25}\log2-\frac{16}{45}\right)\\
&=\frac{2106\log13-468\log2-4160}{925}<\frac{124}{125}.
\end{aligned}
\end{equation}

Together with Step 1, this proves $1-\mathcal R(b)>1/125$ throughout the
required interval.

\textbf{Step 3: bound the gain coefficient.} We now check the other input to \eqref{eq:8.31}.
Clearing the positive denominators in $\delta_0(b)\ge1/12$ gives
\begin{equation}\label{eq:8.37}
6(1+b)(b^2+b+2)-d(1-b)(2+b)(3+b)
=b(26+19b+4b^2-b^3)\ge0.
\end{equation}

This proves both assertions. 
\end{proof}

\textbf{Higher-dimensional part of Theorem~\ref{thm-determinant} completed.} Combining \eqref{eq:8.31} with
Lemma~\ref{lem-rational} gives $\log\det\widehat F\ge\chi/1500$ for $n\ge4$.
The remaining dimensions follow below; the common upper bound and equality
statement are completed in Section~\ref{sec:8.6}.

\textbf{Why dimensions two and three require another estimate.} The scalar formula for
$\mathcal R(b)$ is well defined for $0<b\le1/2$. The preceding bound needs
$\mathcal R(b)<1$; for $n=3$ and $n=2$, however, this scalar remainder exceeds one.
Indeed, on $a\le z\le t$,
\begin{equation}\label{eq:8.38}
\frac1{1-z^2}\ge1+z^2+z^4,\qquad
\frac1{(1-bz)^2}\ge1+2bz+3b^2z^2.
\end{equation}

These follow by multiplying the positive denominators: the two products are
$1-z^6\le1$ and $1-(bz)^3(4-3bz)\le1$. With
$p(b)=4(1-b)^2(2+b)(3+b)/(b^2+b+2)$, exact polynomial integration gives
\begin{equation}\label{eq:8.39}
\begin{aligned}
\mathcal R(b)&\ge p(b)\int_a^t
z^2(1+z^2+z^4)(1+2bz+3b^2z^2)\,dz,\\
\mathcal R(1/3)&>\frac{16}{15}>1,\\
\mathcal R(1/2)&>\frac76>1.
\end{aligned}
\end{equation}

Thus the scalar estimate does not prove a positive determinant deficit in these
two dimensions. We now retain transverse block information, first to settle
dimension two and then to derive the additional control needed in dimension three.

\subsection{Block constraints and dimension two in Theorem~\ref{thm-determinant}}\label{sec:8.3}

The scalar estimate has reached its limitation. To handle the remaining
dimensions, we retain the components of $u_i=s_iv+y_i$ transverse to $v$.
Work in $\mathbb Rv\oplus v^\perp$, with $m=n-1$.
In the normalized certificate, $x$ measures the axial decrease, $E$ is the
transverse perturbation, and $j$ couples the two subspaces.

\begin{lemma}[{block identities and constraints}]\label{lem-moments}
The normalized matrix satisfies
\begin{equation}\label{eq:8.40}
\widehat F=\begin{pmatrix}1-x&j^T\\j&I_m+E\end{pmatrix},
\quad x=\frac{\chi}{2(2+b)},\quad
\sigma:=\operatorname{tr}E=\frac{2\chi}{d(2+b)},\quad x=\frac d4\sigma,
\end{equation}
where
\begin{equation}\label{eq:8.41}
E=-\frac1d\sum_i\nu_i(s_i-b)y_i y_i^T,
\qquad
j=-\frac{\sqrt{1-b}}{b\sqrt d}\sum_i\nu_i(s_i-b)y_i.
\end{equation}

The scalar and matrix constraints are
\begin{equation}\label{eq:8.42}
0\le\sigma\le\frac{2(1-b)}{d+b^2},\qquad 0\le\chi\le2,
\end{equation}
\begin{equation}\label{eq:8.43}
\begin{pmatrix}x&-j^T\\-j&tI_m-E\end{pmatrix}\succeq0,
\qquad -\frac12I_m\prec E\prec tI_m.
\end{equation}

\end{lemma}

The trace of $E$ is nonnegative, but $E$ need not be positive semidefinite.
The Gram constraint limits the coupling: for every transverse vector $q$,
$(q^Tj)^2\le x\,q^T(tI_m-E)q$. At the equality contacts, either $y_i=0$
or $s_i-b=0$, so $E=j=0$.

\begin{proof}
\textbf{Recover the blocks.} The projected John identities are
$\sum_i\lambda_i y_i=\sum_i\lambda_i s_i y_i=0$ and
$\sum_i\lambda_i y_iy_i^T=I_m$. With the reweighted vector and matrix moments
$C_r=\sum_i\nu_i s_i^r y_i$ and $Y_r=\sum_i\nu_i s_i^r y_iy_i^T$, they give
\begin{equation}\label{eq:8.44}
C_0=aC_1,\qquad C_1=aC_2,\qquad Y_0-aY_1=I_m.
\end{equation}

The identity $f(s)=d/2-(s-b)/[2(1-as)]$ gives the transverse expression
for $E$, or equivalently $Y_1-bY_0=-dE$.
Adding the axial term $t/(1-a)$ to \eqref{eq:8.7} and normalizing by $F_0$ gives
$1-x$. Taking the trace of $E$, using $\|y_i\|^2=1-s_i^2$, gives $\sigma$.
The mixed entry follows from $C_1-tC_2$ and \eqref{eq:8.44}.
This proves \eqref{eq:8.40}--\eqref{eq:8.41}.

\textbf{Bound the scalar range.} It suffices to bound $\chi$.
The nonnegative sum $\sum_i\nu_i(1-s_i)(s_i+b)^2$ expands to
$-N_3+(1-2b)N_2+(2b-b^2)N_1+b^2N_0$.
Combining it with \eqref{eq:8.5}--\eqref{eq:8.6} cancels the remaining free moment:
\begin{equation}\label{eq:8.45}
\chi+\frac{(1-b)(2+b)^2}{(d+b^2)^2}
\sum_i\nu_i(1-s_i)(s_i+b)^2
=\frac{d(2+b)(1-b)}{d+b^2}.
\end{equation}

Nonnegativity gives the stated bound on $\sigma$ through \eqref{eq:8.40}.
The estimate $\chi\le2$ also uses
$2(d+b^2)-d(2+b)(1-b)=b(2+3b-b^2)\ge0$.

\textbf{Keep the coupling in a positive block matrix.} The required matrix has
the Gram representation
\begin{equation}\label{eq:8.46}
\begin{pmatrix}x&-j^T\\-j&tI_m-E\end{pmatrix}
=\sum_i\nu_i(1+s_i)r_ir_i^T,\qquad
r_i=\binom{(s_i-b)/\sqrt{2(2+b)}}
{\sqrt{(1-b)(2+b)/(2d)}\,y_i}.
\end{equation}

Its weights are nonnegative, proving the first part of \eqref{eq:8.43}.
Solving the two relations for $Y_0,Y_1$ gives
\begin{equation}\label{eq:8.47}
\begin{aligned}
\sum_i\nu_i(1-s_i)y_i y_i^T&=\frac d{2+b}(I_m+2E),\\
\sum_i\nu_i(1+s_i)y_i y_i^T&=\frac{2d}{(1-b)(2+b)}(tI_m-E).
\end{aligned}
\end{equation}

Both sums are positive definite: the contacts at $s_i=\pm1$ have $y_i=0$,
and all other transverse vectors span $v^\perp$ by the John second moment.
Thus $I_m+2E\succ0$ and $tI_m-E\succ0$, proving the strict spectral bounds. 
\end{proof}

\textbf{Dimension two.} Here $m=1$, so the trace determines the transverse block:
$E=\sigma$, $x=3\sigma/8$, and $t=3/4$.
The remaining coupling loss is $j^2$. Substituting the Gram bound
$j^2\le x(t-\sigma)$ into the determinant gives
\begin{equation}\label{eq:8.48}
\det\widehat F\ge(1-x)(1+\sigma)-x(t-\sigma)
=1+\frac{11}{32}\sigma=1+\frac{11}{60}\chi.
\end{equation}

Dimension three needs further control because the transverse trace alone no
longer determines the transverse determinant. We next state the precise
estimate needed there.

\subsection{A direct bound for transverse deformation and coupling}\label{sec:8.4}

For the two-dimensional transverse block in the next subsection, we need to
control its deformation and its coupling to the axis together. The same
bound will also be used for local stability in Section~\ref{sec:9.3}.
The matrix $K_E=I_m+2E$ is positive definite by Lemma~\ref{lem-moments}.
The nonnegative quantity $A_E=\operatorname{tr}(E^2K_E^{-1})$ measures
transverse deformation: each eigenvalue $e$ of $E$ contributes $e^2/(1+2e)$.

\begin{lemma}[{direct coupling bound}]\label{lem-variational}
Let
$K_E=I_m+2E$, $A_E=\operatorname{tr}(E^2K_E^{-1})$, and
$L=d(1-b)/(d+b^2)-d\sigma/2\ge0$.
If $\sigma>0$, then
\begin{equation}\label{eq:8.49}
A_E\le\frac\sigma2,\qquad
\|K_E^{-1}j\|^2\le L\left(\frac\sigma2-A_E\right).
\end{equation}

\end{lemma}

The first inequality bounds transverse deformation. The second says that
the remaining amount $\sigma/2-A_E$ controls the transformed coupling.
The coefficient $L$ is the distance to the upper trace bound in
\eqref{eq:8.42}, multiplied by $d/2$.

\begin{proof}
\textbf{Evaluate only the tests needed for these conclusions.} We will take $Z=EK_E^{-1}$, first with $q=0$ and then with
$q=K_E^{-1}j/L$ when $L>0$. For these calculations set
\begin{equation}\label{eq:8.50}
\ell=\operatorname{tr}(ZE)+q^Tj,\qquad
Q_2=\operatorname{tr}(Z^2K_E)+4q^TZj+L\|q\|^2.
\end{equation}

To bound this expression directly from the frame, put
$\gamma=\sqrt{d(1-b)}/(d+b^2)$. The mixed moments, transverse identity,
and scalar relation in Section~\ref{sec:8.3} give, respectively,
\begin{equation}\label{eq:8.51}
\begin{aligned}
\sum_i\nu_i(1-s_i)y_iy_i^T&=\frac d{2+b}K_E,\\
\gamma\sum_i\nu_i(1-s_i)(s_i+b)y_i&=\frac{2d}{2+b}j,\\
\gamma^2\sum_i\nu_i(1-s_i)(s_i+b)^2&=\frac d{2+b}L.
\end{aligned}
\end{equation}

Together with the expression for $E$, these identities evaluate the linear
form and its associated squared norm:
\begin{equation}\label{eq:8.52}
\begin{aligned}
\ell&=-\frac1d\sum_i\nu_i(s_i-b)y_i^T
\bigl(Zy_i+\gamma(s_i+b)q\bigr),\\
\sum_i\nu_i(1-s_i)\|Zy_i+\gamma(s_i+b)q\|^2
&=\frac d{2+b}Q_2.
\end{aligned}
\end{equation}

Weighted Cauchy--Schwarz, with
$\|y_i\|^2=(1-s_i)(1+s_i)$, therefore gives
$\ell^2\le\chi Q_2/[d(2+b)]=(\sigma/2)Q_2$.
Pole terms have $y_i=0$, so this calculation requires no division by
$1\pm s_i$. Completing one scalar square yields
\begin{equation}\label{eq:8.53}
2\ell-Q_2\le2\ell-\frac{2\ell^2}{\sigma}
=\frac\sigma2-\frac2\sigma\left(\ell-\frac\sigma2\right)^2
\le\frac\sigma2.
\end{equation}

\textbf{Extract the deformation and coupling bounds.}
The chosen $Z=EK_E^{-1}$ is symmetric and makes
$\operatorname{tr}(ZE)=\operatorname{tr}(Z^2K_E)=A_E$.
With $q=0$, this gives $2\ell-Q_2=A_E$, proving the first bound.
Keeping this $Z$, use $2EK_E^{-1}=I_m-K_E^{-1}$: the part depending on $q$
is $2q^TK_E^{-1}j-L\|q\|^2$. For $L>0$ it is maximized at the chosen
$q=K_E^{-1}j/L$, giving
\begin{equation}\label{eq:8.54}
2\ell-Q_2=A_E+\frac1L\|K_E^{-1}j\|^2\le\frac\sigma2.
\end{equation}

If $L=0$, the same finite Cauchy--Schwarz calculation with $Z=0$ and $q=j$
gives $\|j\|^4\le0$, hence $j=0$. This proves the second bound at that endpoint.
The case $\sigma=0$ has $E=j=0$ by \eqref{eq:8.40}--\eqref{eq:8.41}, so the estimates extend
there as well. 
\end{proof}

The useful output is \eqref{eq:8.49}. Section~\ref{sec:10.4} records the more general inequality
for arbitrary tests; the proof here has already supplied everything needed
for the three-dimensional and local stability arguments.

\subsection{Dimension three in Theorem~\ref{thm-determinant}: a single polynomial certificate}\label{sec:8.5}

The transverse block is now two-dimensional, so its trace no longer determines
its determinant. We must control both the transverse determinant and its
interaction with $j$. The proof first expresses that interaction through two
scalar quadratic forms, then combines the Gram and variational inequalities
to remove them. The remaining task is a polynomial bound on an interval.

\begin{proposition}[{three-dimensional bound}]\label{prop-dim3}
In dimension three,
\begin{equation}\label{eq:8.55}
\det\widehat F\ge1+\frac2{11}\sigma
=1+\frac{36}{385}\chi.
\end{equation}

\end{proposition}

\begin{proof}
The desired output is a positive determinant excess proportional
to $\sigma$. The coefficient $2/11$ below is the margin supplied by the
polynomial estimate.

\textbf{Step 1: describe the transverse block by scalar invariants.} In dimension
three, the definitions specialize to $E\in\mathbb S^2$, $x=5\sigma/12$,
$t=2/3$, and $0\le\sigma\le3/4$. If $\sigma=0$, the defect vanishes and
\eqref{eq:8.41} gives $\widehat F=I$. Suppose $\sigma>0$ and write
\begin{equation}\label{eq:8.56}
p=\det E,\qquad z=\sigma+4p,\qquad
L=\frac58-\frac56\sigma,\qquad
\Delta=\frac{11}{3}(1+\sigma)-z.
\end{equation}

Here $p$ is the second invariant of $E$ besides its trace. The combination
$z=\sigma+4p$ is chosen because it expresses the variational remainder
$\sigma-2A_E$ by a simple fraction. The coefficient $L$ is exactly the one
in Lemma~\ref{lem-variational} at $n=3$. The quantity $\Delta$ will be the positive denominator
when we combine the two coupling inequalities.
For a two-by-two matrix, Cayley--Hamilton gives
\begin{equation}\label{eq:8.57}
\det K_E=1+\sigma+z,\qquad
\sigma-2A_E=\frac{z}{1+\sigma+z}.
\end{equation}

Lemma~\ref{lem-variational} and $\det E\le\sigma^2/4$ therefore imply
\begin{equation}\label{eq:8.58}
0\le z\le\sigma(1+\sigma),\qquad
j^TK_E^{-2}j\le\frac{Lz}{2(1+\sigma+z)},\qquad \Delta>0.
\end{equation}

Thus $z$ lies in an explicit interval once $\sigma$ is fixed. This will later
make a concavity argument possible.

\textbf{Step 2: bound exactly the coupling term in the determinant.} For a
two-by-two transverse block, that term is
$j^T\operatorname{adj}(I_2+E)j=(1+\sigma)\|j\|^2-j^TEj$.
Use the short names $r=\|j\|^2$ and $\omega=j^TEj$ for these two quadratic
forms, and set $N=tI_2-E\succ0$. The Gram constraint and \eqref{eq:8.58} read, respectively,
\begin{equation}\label{eq:8.59}
(t-\sigma)r+\omega\le x\det N,
\end{equation}
\begin{equation}\label{eq:8.60}
\bigl[(1+\sigma)(1+4\sigma)-z\bigr]r-4(1+\sigma)\omega
\le\frac L2z(1+\sigma+z).
\end{equation}

Both inequalities are linear in $r,\omega$. Choose a positive combination
whose coefficients are exactly those of $(1+\sigma)r-\omega$.
Solving this two-variable coefficient-matching problem gives the multipliers
$[3(1+\sigma)+z]/\Delta$ and $5/(3\Delta)$, respectively.
Their positivity preserves the inequality direction, and their common
denominator explains the definition of $\Delta$. Adding yields
\begin{equation}\label{eq:8.61}
j^T\operatorname{adj}(I_2+E)j
\le\frac{x\det N[3(1+\sigma)+z]+\frac56Lz(1+\sigma+z)}\Delta.
\end{equation}

\textbf{Step 3: eliminate the matrix variables.} Substitute this bound into
$\det\widehat F=(1-x)(1+\sigma+p)-j^T\operatorname{adj}(I_2+E)j$
and use $p=(z-\sigma)/4$. Multiplication by $108\Delta$ clears the variable
denominator and gives a polynomial in $\sigma,z$ alone; $108$ is a convenient
positive normalization of its rational coefficients:
\begin{equation}\label{eq:8.62}
108\Delta(\det\widehat F-1)\ge\mathcal P_*(\sigma,z),
\end{equation}
\begin{equation}\label{eq:8.63}
\mathcal P_*(\sigma,z)=72\sigma(1+\sigma)
+\left(\frac{171}{4}-\frac{53}{4}\sigma+75\sigma^2\right)z
+\left(-\frac{333}{4}+75\sigma\right)z^2.
\end{equation}

\textbf{Step 4: use concavity to check only the endpoints.} The coefficient of
$z^2$ is negative for $0\le\sigma\le3/4$, so the minimum over the admissible
$z$ interval occurs at an endpoint. At $z=0$ the polynomial is
$72\sigma(1+\sigma)$. We show that the other endpoint is at least this large.
At $z=\sigma(1+\sigma)$, it equals
\begin{equation}\label{eq:8.64}
\frac{\sigma(1+\sigma)}4
\bigl(300\sigma^3+267\sigma^2-386\sigma+459\bigr)
\ge72\sigma(1+\sigma),
\end{equation}
because the cubic in parentheses minus $288$ has the nonnegative decomposition
\begin{equation}\label{eq:8.65}
300\sigma^3+267\left(\sigma-\frac34\right)^2
+\frac{29}{2}\sigma+\frac{333}{16}.
\end{equation}

Consequently $\mathcal P_*\ge72\sigma(1+\sigma)$ throughout the admissible
interval. Returning through \eqref{eq:8.62}, the bound
$\Delta\le11(1+\sigma)/3$ cancels $1+\sigma$ and gives
$\det\widehat F-1\ge2\sigma/11$.
Thus $2/11=72/[108(11/3)]$ is the final margin left by the polynomial bound. 
\end{proof}

Keeping $L=5/8-5\sigma/6$ is decisive: a larger coefficient weakens the second constraint and loses the single-interval polynomial certificate. The proof uses only the moment identities, positive-semidefinite blocks, two-by-two matrix identities, and polynomial inequalities with rational coefficients.

\subsection{Completion of Theorems~\ref{thm-determinant} and \ref{thm-frame-bound}}\label{sec:8.6}

The three branches now supply determinant excess. We put them on the same
logarithmic scale, then recover the upper bound and the equality condition.

\textbf{Obtain one lower bound for all dimensions.}
For $n\ge4$, Lemma~\ref{lem-rational} and \eqref{eq:8.31} give $\log\det\widehat F\ge\chi/1500$.
For $n=2,3$, equations \eqref{eq:8.48} and \eqref{eq:8.55} give $\det\widehat F\ge1+c\chi$ with $c=11/60$ and $c=36/385$, respectively. Since $\chi\le2$ and $\log(1+u)\ge u/(1+u)$ for $u\ge0$,
\begin{equation}\label{eq:8.66}
\log\det\widehat F\ge\frac{c}{1+2c}\chi\ge\frac{\chi}{1500}.
\end{equation}

\textbf{Use the trace for the upper bound.} This part is common to all dimensions:
\begin{equation}\label{eq:8.67}
\log\det\widehat F\le\operatorname{tr}\widehat F-n
=\sigma-x=\frac{\chi}{2d}.
\end{equation}

\textbf{Recover equality from zero defect.} The lower bound forces $\chi=0$ when
$\det\widehat F=1$. Every summand in \eqref{eq:8.3} is nonnegative with a positive weight
and denominator, so this is equivalent to $s_i\in\{-1,b\}$ for all $i$.
Conversely, this contact condition gives $x=j=E=0$ by \eqref{eq:8.40}--\eqref{eq:8.41}.
This completes the proof of Theorem~\ref{thm-determinant} announced at the start of Section~\ref{sec:8}.

By \eqref{eq:8.2}, this proves \eqref{eq:3.12}. At equality $F=F_0$, hence $Q=M_0$; together with the positivity established above, this completes Theorem~\ref{thm-frame-bound}.
